% --- Ajustes Finos do Sumário (Normas UEMA) ---

% 1. Remove TODOS os espaçamentos verticais extras (espaço 1.5 uniforme)
\setlength{\cftbeforechapterskip}{0pt}
\renewcommand{\insertchapterspace}{}

% 2. Deixa os pontinhos da linha mais juntos
\renewcommand{\cftdotsep}{2}
\renewcommand{\cftchapterdotsep}{2}
\renewcommand{\cftsectiondotsep}{2}
\renewcommand{\cftsubsectiondotsep}{2}
\renewcommand{\cftsubsubsectiondotsep}{2}

% 3. Alinha TODAS as seções e subseções rigorosamente à esquerda (margem 0)
\setlength{\cftchapterindent}{0pt}
\setlength{\cftsectionindent}{0pt}
\setlength{\cftsubsectionindent}{0pt}
\setlength{\cftsubsubsectionindent}{0pt}

% 4. Distância entre o número e o texto
\setlength{\cftchapternumwidth}{1.3cm}
\setlength{\cftsectionnumwidth}{1.3cm}
\setlength{\cftsubsectionnumwidth}{1.3cm}
\setlength{\cftsubsubsectionnumwidth}{1.3cm}

% 5. Força todas as subseções a ficarem em negrito
\renewcommand{\cftsectionfont}{\bfseries}
\renewcommand{\cftsubsectionfont}{\bfseries}
\renewcommand{\cftsubsubsectionfont}{\bfseries}

% 5. Números das páginas nunca em negrito
\renewcommand{\cftchapterpagefont}{\normalfont}
\renewcommand{\cftsectionpagefont}{\normalfont}
\renewcommand{\cftsubsectionpagefont}{\normalfont}
\renewcommand{\cftsubsubsectionpagefont}{\normalfont}